\documentclass[11pt,a4paper]{article}
\usepackage[margin=25mm]{geometry}
\usepackage{amsmath,amssymb,amsthm}
\newtheorem{lemma}{Lemma}\newtheorem{theorem}{Theorem}
\title{Irreducible parabolic kernels and the pair count of type $(6,7)$}
\author{Human Compression project}\date{}
\begin{document}\maketitle
Let $G=O'N$. Write $x_J(p,q)$ for the number of ordered pairs
$(r,s)$ in $J$ with $|r|=p$, $|s|=q$, $|rs|=2$, and write $e_J(p,q)$
for the number which generate $J$. We recover the fixed
Connor--Leemans Table~1 entry
\[                         e_G(6,7)/(2|G|)=354.                \tag{1}\]
The subgroup roster, primary class data, exact character evaluations
and accepted projective models are borrowed premises.

\begin{lemma}\label{parabolic}
A pair of type $(6,7)$ in a row 9 or 10 parabolic $P$ either generates
$P$ or is contained in a maximal subgroup of row 7.
\end{lemma}
\begin{proof}
In the accepted projective model, $P=V:Q$, where $V\cong C_7^2$ and
$Q\cong\mathrm{GL}_2(7)/C_3$. The action on $V$ is irreducible:
its restriction to the determinant-one block is the natural
irreducible two-dimensional action, or its dual. The scalar quotient
does not change this assertion.

Let $J$ be a maximal proper subgroup of $P$. If $V$ is not contained
in $J$, maximality gives $JV=P$. Then $J\cap V$ is invariant under
$J$ and, since $V$ is abelian, under $V$; hence it is $P$-invariant.
Irreducibility forces $J\cap V=1$. Thus $J$ is a complement of
order 672. The complete maximal lists identify these complements
as row 115 or 116. Both classes occur below row 7, so every pair
inside a complement has a row 7 overgroup.

All other maximal subgroups contain $V$. Their quotient orders,
from the complete lists, are 336, 84, 32 and 24. The quotient of
order 336 is the unique index-two subgroup
$\mathrm{SL}_2(7)$ of $Q$: it is perfect, of index two, and hence
equals $Q'$. It has only the central involution $-I$.
Projection modulo $V$ preserves the exact orders six and two,
since those orders are coprime to seven. The order-seven entry has
image of order one or seven. Such factors cannot multiply to $-I$:
if $\bar r\bar s=-I$, then $\bar s=-\bar r^{-1}$ has order dividing
six, and is not the identity because $\bar r$ has order six.

A group of order 84 has a normal Sylow 7-subgroup: its number divides
12 and is congruent to one modulo seven. Modulo that subgroup the
order-six factor retains order six and the order-seven factor is
killed, contradicting a product of order two. Finally, in quotients
of orders 32 or 24 the order-seven factor is already killed, giving
the same contradiction. Thus none of these four maximal types
contains a pair of type $(6,7)$. A proper generated subgroup lies
in a maximal subgroup, so only the complement case remains.
\end{proof}

\begin{lemma}\label{outside}
Every proper subgroup generated by a pair of type $(6,7)$ and
contained in no row 7 maximal has a unique maximal overgroup in $G$.
\end{lemma}
\begin{proof}
The flag normalizer of row 22 has a normal Sylow 7-subgroup. Modulo
it, an order-six element retains order six and an order-seven
element is killed. It therefore has no pair of the required type.
Using the cyclic-six and cyclic-seven ancestor sets, remove $G$,
the descendants of row 7 and row 22, and proper descendants of the
two parabolic classes by Lemma~\ref{parabolic}. The remaining
possible generated classes are exactly
\[
                 2,3,4,5,6,9,10,21,43,44,68,172,174,237.
\]
Classes 2, 3, 6, 21, 43 and 44 are maximal. Classes 4, 5, 9, 10,
172 and 174 have unique actual maximal parents, from the accepted
direct incidence calculations along
\[
                 9\longrightarrow4\longrightarrow2,
 \qquad10\longrightarrow5\longrightarrow3,
 \qquad172\longrightarrow4,
 \qquad174\longrightarrow5.
\]

It remains to exclude 68 and 237. The accepted row 21 model has
normal kernel $C_4^3$ and quotient $L_3(2)$. The preimage of the
maximal subgroup $7:3$ is row 68, the only maximal class of order
1344 in row 21. Thus row 68 has a normal 2-subgroup and quotient
of order 21. In row 237 the unique subgroup of order 56 is normal; its
characteristic elementary abelian Sylow 2-subgroup of order eight
is therefore normal in row 237, again with quotient of order 21.
In either quotient, an order-six element has image of exact order
three: its order-three subgroup cannot meet the 2-kernel. An
order-seven element retains order seven. Their product, being an
involution before projection, has trivial image in this odd-order
quotient. This would make the two images inverse, contradicting
their different orders. Neither class contains the pair. Every
remaining class has the required unique maximal overgroup.
\end{proof}

\begin{lemma}\label{common}
A pair of type $(6,7)$ has at most one common commuting involution.
If it lies both in a row 7 maximal $C_G(t)$ and in a projective
maximal $M$, then $t\in M$. No other maximal class with nonzero
pair count can meet row 7 in such a pair.
\end{lemma}
\begin{proof}
The primary centralizer orders are 72 for the order-six class, and
1372 or 49 for the order-seven classes. Thus the common centralizer
has order dividing four. In the case 1372, its Sylow 2-subgroup is
cyclic of order four, because the primary power data give an
order-28 root. The common centralizer is cyclic and has at most
one involution.

The possible generated classes common to row 7 and a projective
maximal, after excluding row 22, are exactly
\[
65,66,110,111,112,113,115,116,117,118,171,173,175,176,177.
\]
For all but 177, the accepted $(7,8)$ common-involution argument
applies word for word to these classes. The block subgroups contain
the central involution of their perfect $\mathrm{SL}_2(7)$ subgroup;
the outer-centralizer classes have their direct central involution;
and the centerless classes 171 and 173 have their normalizers inside
each projective maximal, by their unique simple overgroups.
Only the uniqueness of the common involution depends on the pair
type, and that was just proved. Class 177 is the index-two subgroup
$\mathrm{PSL}_2(7)\times C_2$ in a row 110 or 111 outer centralizer;
the complete index-two list identifies it, the other two being
centerless copies of $\mathrm{PGL}_2(7)$. Its central involution
belongs to every overgroup and must equal $t$.

The only other maximal classes with nonzero $x_M(6,7)$ are 6, 21,
43 and 44. The cyclic-seven ancestor set has empty intersection
with the common descendants of any of these and row 7. This
excludes a common pair and completes the proof.
\end{proof}

\begin{theorem}
Let $H$ and $K$ be the accepted inner and outer involution centralizers
in $L_3(7):2$, of orders 1344 and 672. Then
\[
 e_G(6,7)=x_G(6,7)-\sum_M[G:M]x_M(6,7)
       +2\cdot342868680\bigl(x_H(6,7)+2x_K(6,7)\bigr),          \tag{2}
\]
where $M$ runs through the thirteen maximal conjugacy classes.
The ordered automorphism-orbit count is 354.
\end{theorem}
\begin{proof}
By Lemma~\ref{outside}, a nongenerating pair outside all row 7
maximals appears exactly once in the maximal sum. By
Lemma~\ref{common}, a pair in row 7 appears once there; its excess
is the number of its projective maximal overgroups, and within
each of them the defining involution belongs to that maximal.
Each projective class has 122760 members, with 2793 inner and
5586 outer involutions per member. Their inner incidence total is
342868680, giving (2), with no repeated pair in a fixed maximal.
The retained literal certificate gives
\[
                        x_H(6,7)=1344,\qquad x_K(6,7)=1008.
\]
The character evaluation gives $x_G(6,7)=380748811766400$ and
maximal sum $56795511104640$. The correction is $2304077529600$.
Substitution and division by $|\operatorname{Aut}(G)|=921631011840$
give (1). The action on ordered generating pairs is free; there
is no quotient by interchange of the entries.
\end{proof}

All finite values come from the retained order-seven calculation.
The independent pointwise reader verifies its complete class and
partner lists. The new contribution is the irreducible-kernel
argument, which removes every proper parabolic case at once,
and its source-bound correction of the maximal sum. The original
numerical assertion is due to Connor and Leemans. No chirality,
whole-source, open-question or Lean credit is claimed.
\end{document}
